Five pairings. On covector pairs both sides are $G$. Against $\nO$: $h_u(\alpha, \nO + \nu\tilde v) = \nu\alpha(u) + \nu G(\alpha, \tilde v) = \nu\alpha(u + v) = h_{u+v}(\alpha, \nO)$, identically with $\mu$ against $\nInf$. On $\nO$ with itself: $2\nu^2\tilde v(u) + \nu^2 G(\tilde v, \tilde v) = \nu^2(G(\tilde v, \tilde v) + 2\tilde v(u)) = 0$ by the normalisation, and identically with $\mu^2$; on the cross pairing the same combination appears with coefficient $\mu\nu$ and cancels, leaving $-1$. Unipotence: $(\Phi_v - \id)$ has image in $T^*M$ and square zero. An isometry of quadratic modules induces an isomorphism of Clifford algebras, grading-preserving since it is induced by a map of the generating module; it is the identity on $T^*M$, hence on the subalgebra $T^*M$ generates. The cocycle: both sides are the identity on $T^*M$ and send $\nO$ to $\nO + \nu(\tilde v + \tilde w)$, $\nInf$ to $\nInf + \mu(\tilde v + \tilde w)$, with $\tilde v$ normalised against $u$ and $\tilde w$ against $u + v$; so it suffices that $\tilde v + \tilde w$ is the representative of $v + w$ normalised against $u$. It represents $v + w$ by linearity of the solve, and $G(\tilde v + \tilde w, \tilde v + \tilde w) + 2(\tilde v + \tilde w)(u) = [G(\tilde v, \tilde v) + 2\tilde v(u)] + [G(\tilde w, \tilde w) + 2\tilde w(u) + 2G(\tilde w, \tilde v)] = 0 + [G(\tilde w, \tilde w) + 2\tilde w(u + v)] = 0$, using $G(\tilde w, \tilde v) = \tilde w(v)$, which is the defining property of $\tilde v$. The images: $\Phi_v(k^*) = \Phi_v(\nO)/2\nu = k^* + \tfrac12\tilde v$, and since $\tilde v k + k\tilde v = 2G(\tilde v, k) = 0$, $\Phi_v(\sigma) = (k^* + \tfrac12\tilde v)k - k(k^* + \tfrac12\tilde v) = \sigma + \tfrac12(\tilde v k - k\tilde v) = \sigma + \tilde v k$.
